\chapter{The workspace}

\label{ch:workspace}

This chapter answers where a behaviour lives. pii is a Cargo workspace of twenty-three members ---
twenty-two crates and \texttt{xtask} --- and the split is by \emph{responsibility that changes for
its own reasons}, not by layer. The practical consequence for a reader is that a symptom names its
own crate: a refusal raised while a graph is being read comes from \texttt{pii-graph}; a run that
never finishes comes from \texttt{pii-server}'s engine; a duplicated implementation is found by
\texttt{xtask}, not by a reviewer. The chapter ends with a table that maps a symptom to the file a
reader should open first, because that is the question a new contributor actually asks.

Two rules keep the workspace from decaying, and both are enforced mechanically rather than by
convention: a file has a budget of 450 effective lines, and a copy of an idea is treated as a
defect even when both copies are correct today. The second rule is the unusual one, and it is why
several crates contain a file whose only job is to be the single home of a concept.

\section{Membership}

\begin{longtable}{@{}p{0.17\textwidth}p{0.44\textwidth}p{0.31\textwidth}@{}}
	\toprule
	\textbf{Crate} & \textbf{What it owns} & \textbf{A file you will open} \\
	\midrule
	\endhead
	\texttt{pii-abi} & The vocabulary: ids, ports, run status, journal events, capability bits, the
	wire types shared by everything. It depends on no other pii crate. & \texttt{src/journal.rs},
	\texttt{src/cap.rs} \\
	\texttt{pii-core} & The few root concepts a node needs to run: the context, the capability set,
	the filesystem scope, the budget, the clock, the state directory resolver. & \texttt{src/ctx.rs},
	\texttt{src/cap.rs} \\
	\texttt{pii-observe} & The closed event vocabulary and the sinks that receive it. Names live in
	one module so a metric cannot be spelled two ways. & \texttt{src/names.rs} \\
	\texttt{pii-store} & Every redb table and its write paths: runs, queue, approvals, heals, memory,
	skills. Deliberately knows nothing about graphs. & \texttt{src/queue.rs} \\
	\texttt{pii-graph} & The graph spec, the loader, static validation, the splice, the scheduler,
	the journal writer, the tick engine. & \texttt{src/fragment/splice.rs}, \texttt{src/sched/} \\
	\texttt{pii-server} & The engine bridge, the daemon, the node kinds, the transports, the JSON-RPC
	dispatch table. & \texttt{src/engine/runs.rs} \\
	\texttt{pii-proto} & The client/server protocol and its method names. & \texttt{src/run.rs} \\
	\texttt{pii-uds} & The unix socket the CLI and the daemon speak over. & --- \\
	\texttt{pii-cli} & The \texttt{pii} binary: argument parsing and one module per command group.
	& \texttt{src/cmd\_run.rs} \\
	\texttt{pii-worker} & The agent worker: argv construction, approval modes, in-flight limits.
	& --- \\
	\texttt{pii-model} & Model providers and the model id vocabulary. & --- \\
	\texttt{pii-net} & The network surface: egress grants, dial and accept. & --- \\
	\texttt{pii-memory} & Entities, tiers, the hub the memory nodes read. & --- \\
	\texttt{pii-learn} & Lessons, candidates, promotions, incidents. & --- \\
	\texttt{pii-tools} & The tool surface an agent may call, and its sandbox. & --- \\
	\texttt{pii-channels} & Outbound channels: Discord, MQTT. & --- \\
	\texttt{pii-tui}, \texttt{pii-web}, \texttt{pii-mcp} & Three operator surfaces over the same
	protocol. & --- \\
	\texttt{pii-wasm}, \texttt{pii-hal}, \texttt{pii-evolve} & Optional capability surfaces, each
	behind its own feature. & --- \\
	\bottomrule
\end{longtable}

The table is worth reading twice, because the absences are as deliberate as the presences.
\texttt{pii-store} has no graph types in it: a table stores an event, not a \texttt{RunState}. That
is what allows the store to be opened by a lint registry, a test fixture or a repair tool without
dragging the scheduler's semantics along. Conversely \texttt{pii-graph} does not contain a node
kind: the built-in seven (\texttt{transform}, \texttt{route}, \texttt{emit}, \texttt{wait},
\texttt{timer}, \texttt{refuse}, \texttt{require}) live beside the registry, and everything that
touches a model, a container, a store or a channel is registered by \texttt{pii-server}. A crate
that could run a graph on its own would need none of those, which is what makes the graph crate
testable with a stub registry.

\section{The dependencies that are allowed}

\begin{center}
\resizebox{\linewidth}{!}{%
\begin{tikzpicture}[node distance=8mm and 12mm]
	\node[b] (abi) {\texttt{pii-abi}\\\scriptsize vocabulary};
	\node[b, right=of abi] (core) {\texttt{pii-core}\\\scriptsize ctx, caps, budget};
	\node[b, right=of core] (graph) {\texttt{pii-graph}\\\scriptsize spec, splice, sched};
	\node[b, right=of graph] (server) {\texttt{pii-server}\\\scriptsize engine, kinds, daemon};
	\node[b, right=of server] (cli) {\texttt{pii-cli}\\\scriptsize the binary};
	\node[s, below=of graph] (store) {\texttt{pii-store}\\\scriptsize redb tables};
	\node[b, below=of core] (obs) {\texttt{pii-observe}\\\scriptsize events};
	\node[m, below=of server] (model) {\texttt{pii-model}\\\scriptsize providers};

	\draw[f] (abi) -- (core);
	\draw[f] (core) -- (graph);
	\draw[f] (graph) -- (server);
	\draw[f] (server) -- (cli);
	\draw[f] (graph) -- (store);
	\draw[f] (abi) -- (obs);
	\draw[f] (core) -- (obs);
	\draw[f] (server) -- (model);
	\draw[sf] (store) to[bend right=16] node[p,below=2pt,midway]{never} (graph);
	\draw[sf] (cli) to[bend left=12] node[p,above=2pt,midway]{in process} (server);
\end{tikzpicture}}
\captionof{figure}{The dependency direction, left to right; the dashed edge marked \enquote{never}
is the rule that keeps the store reusable: no table may name a graph type; the CLI's edge is worth
noticing because it is why an operator running \texttt{pii run start} needs no daemon.}
\end{center}

Read from each crate's own manifest, the edges are these:

\begin{center}
\footnotesize
\begin{tabular}{@{}p{0.16\textwidth}p{0.76\textwidth}@{}}
	\toprule
	\textbf{Crate} & \textbf{Depends on (pii crates only)} \\
	\midrule
	\texttt{pii-abi} & none \\
	\texttt{pii-observe} & none \\
	\texttt{pii-core} & abi, observe \\
	\texttt{pii-store} & abi, observe \\
	\texttt{pii-proto} & observe \\
	\texttt{pii-graph} & abi, core, observe, store \\
	\texttt{pii-model} & abi, core, observe, net \\
	\texttt{pii-tools} & abi, channels, core, model, observe, proto, uds \\
	\texttt{pii-server} & abi, core, store, observe, graph, proto, model, memory, hal, wasm,
	learn, tools, net, worker, channels \\
	\texttt{pii-cli} & abi, core, store, graph, observe, learn, memory, proto, server, uds,
	worker, tools, tui \\
	\bottomrule
\end{tabular}
\end{center}

Two of these edges are worth a sentence each. \texttt{pii-graph} depends on \texttt{pii-store},
which looks like an inversion until you see what it uses it for: the journal writer takes a
\texttt{Store} and appends events in one transaction, so the crate that decides \emph{what} to
record is also the crate that records it, and the ordering between a run's header and its events is
decided in one place. The reverse edge is what must never exist, and does not: \texttt{pii-store}'s
manifest lists two pii crates, neither of them the graph.

The second edge is \texttt{pii-cli} to \texttt{pii-server}. The binary embeds the daemon in-process
for the commands that need one, which is why \texttt{pii run start} works in a fresh shell with no
socket in sight, and why an operator debugging a stuck run can attach to the same process that
accepted it. The cost is that the CLI links the largest crate in the workspace; the benefit is that
there is one code path for \enquote{start a run}, not a client path and a server path that can
drift.

\section{One behaviour, followed across the crates}

Abstractions about layering are easier to trust with a concrete trace. Take the refusal this manual
returns to repeatedly: a host graph calls a callee whose capability list is wider than the call
grants. The operator sees one line. The line is produced four crates away from where the decision
was taken.

\begin{center}
\resizebox{\linewidth}{!}{%
\begin{tikzpicture}[node distance=8mm and 12mm]
	\node[b] (file) {\texttt{graphs/queue.toml}\\\scriptsize \texttt{caps = [spawn, proc, ...]}};
	\node[b, right=of file] (load) {\texttt{pii-graph}\\\scriptsize \texttt{expand\_one}};
	\node[v, right=of load] (widen) {excess?};
	\node[b, right=of widen] (err) {\texttt{PiiError::Invalid}\\\scriptsize \texttt{graph.call.widen}};
	\node[b, below=of err] (cli) {\texttt{pii-cli}\\\scriptsize prints the message};
	\node[b, below=of load] (server) {\texttt{pii-server}\\\scriptsize \texttt{spawn\_run}};
	\node[s, below=of file] (journal) {journal\\\scriptsize nothing written};

	\draw[f] (file) -- (load);
	\draw[f] (load) -- (widen);
	\draw[f] (widen) -- node[p,above=2pt,midway]{yes} (err);
	\draw[f] (err) -- (cli);
	\draw[sf] (load) -- (server);
	\draw[f] (widen.south) to[bend left=12] node[p,below=2pt,midway]{no} (server);
	\draw[f] (server) -- (journal);
\end{tikzpicture}}
\captionof{figure}{A load-time refusal; the graph is refused before a run exists, so the store stays
untouched and the message an operator reads is the loader's own text --- not a status derived from
a run that failed later.}
\end{center}

The trace in words. \texttt{pii-server}'s \texttt{spawn\_run} is given a \texttt{GraphSpec}, which
means the text has already been read and expanded by \texttt{pii-graph}. Expansion happens in
\texttt{fragment/call.rs}: for each \texttt{call} node it reads the callee, checks the pin against
the bytes, checks the callee's node capabilities against the caps the call grants, and refuses on
the first excess. The error is a \texttt{PiiError::Invalid} whose message begins
\texttt{graph.call.widen} and names the node and the capability. \texttt{pii-server} does not
inspect it; \texttt{pii-cli} prints it. And because the refusal happened during load, no run header
was written: the operator sees a failed start, not a failed run, and there is no half-run to clean
up.

The direction of that flow is the point of the crate split. A decision taken in the loader is not
re-litigated by the engine, and the engine has no code path that could accept a graph the loader
rejected. Every caller that runs or checks a graph goes through the same entry
(\texttt{load\_host\_with\_graphs}), so \enquote{check passed} and \enquote{run started} cannot
disagree.

\section{Inside \texttt{pii-server}}

The largest crate is organised by concern, and three directories carry almost all of the behaviour
this book describes.

\begin{center}
\resizebox{\linewidth}{!}{%
\begin{tikzpicture}[node distance=8mm and 12mm]
	\node[g, align=center, minimum width=3.7cm, minimum height=1.7cm] (eng)
		{\texttt{engine/}\\[2pt]\scriptsize runs $\cdot$ dispatch $\cdot$ orphans\\\scriptsize
		registry $\cdot$ run\_scope};
	\node[g, align=center, minimum width=3.4cm, minimum height=1.7cm, right=12mm of eng] (kinds)
		{\texttt{node\_kinds/}\\[2pt]\scriptsize agent $\cdot$ nodes\\\scriptsize
		compaction};
	\node[g, align=center, minimum width=3.2cm, minimum height=1.7cm, right=12mm of kinds] (daemon)
		{\texttt{daemon/}\\[2pt]\scriptsize config $\cdot$ scheduler\\\scriptsize lag};
	\node[b, below=of eng] (graph) {\texttt{pii-graph}\\\scriptsize the spec it runs};
	\node[b, below=of kinds] (store) {\texttt{pii-store}\\\scriptsize the journal it writes};
	\node[b, below=of daemon] (socket) {the socket\\\scriptsize the RPC table};

	\draw[f] (graph) -- (eng);
	\draw[f] (eng) -- (store);
	\draw[f] (eng) -- (kinds);
	\draw[f] (kinds) -- (daemon);
	\draw[f] (socket) -- (daemon);
\end{tikzpicture}}
\captionof{figure}{The three directories inside \texttt{pii-server} this manual spends its time in: \texttt{engine/}
bridges a graph and the outside world, \texttt{node\_kinds/} is the vocabulary, and \texttt{daemon/}
owns the process.}
\end{center}

\texttt{engine/} is the bridge. \texttt{runs.rs} accepts a graph and returns immediately after the
header and the \texttt{Started} event are durable; \texttt{dispatch.rs} is the JSON-RPC table the
socket reads; \texttt{orphans.rs} is boot recovery for agents a dead daemon left running;
\texttt{run\_scope.rs} builds the capability set and the filesystem scope for a run;
\texttt{registry.rs} assembles the node registry, including a lint-only variant whose stores are in
memory; and \texttt{refusal.rs} records a quarantine refusal as a decision rather than as a log
line.

\texttt{node\_kinds/} is deliberately not a flat directory of thirty files. \texttt{agent/} holds
the turn loop and its tools; \texttt{nodes/} holds the kinds that touch the world, and
\texttt{nodes/coder/} holds the parts of the coder that were extracted when the file crossed the
line budget. \texttt{compaction.rs} is not a node kind at all --- it is agent support used by the
turn loop --- and lives here because that is the only caller.

The daemon owns the process rather than the work: \texttt{config.rs} parses one TOML file with
\texttt{deny\_unknown\_fields} on every section, \texttt{lag.rs} recovers runs a lagging bus left
behind, and the scheduler owns the timer wheel. A run outliving a restart is not the daemon's
problem; it is the journal's, which is why recovery lives in \texttt{engine/orphans.rs} and in
\texttt{pii-graph}'s replay rather than in the daemon's start-up code.

\section{Inside \texttt{pii-graph}}

The graph crate is where a file becomes a runnable object, and its module names are the chapter
titles of the next three chapters.

\begin{itemize}
	\item \texttt{spec.rs} --- the types: \texttt{GraphSpec}, \texttt{NodeSpec}, \texttt{Edge},
	\texttt{EdgeKind}, \texttt{JoinPolicy}, \texttt{GraphLimits}.
	\item \texttt{toml\_load.rs} --- text to types, with the grant ceiling applied as the graph is
	built.
	\item \texttt{validate.rs} --- the static rules, \texttt{V0} to \texttt{V11}, each with its own
	message.
	\item \texttt{fragment/} --- inclusion: \texttt{call.rs} expands \texttt{call} nodes,
	\texttt{splice.rs} weaves \texttt{[[attach]]} fragments, \texttt{call\_output.rs} rewrites host
	values and wires refusals back.
	\item \texttt{sched/} --- readiness, edge states, dead propagation, dispatch of one node.
	\item \texttt{run/} --- the tick loop, resume by replay, retry of a single node.
	\item \texttt{journal.rs} --- the write-behind journal writer and its sequence numbers.
\end{itemize}

The order of that list is the order of a load, and it is worth keeping in mind when a message seems
to come from the wrong place. A port error is raised by \texttt{validate.rs} after the splice, not
by the loader, which is why \texttt{V1} names the node and not the file.

\section{The two budgets}

The first budget is arithmetic. \texttt{xtask/src/lint\_files.rs} defines
\texttt{MAX\_EFFECTIVE\_LINES = 450}, where \enquote{effective} excludes blank and comment lines,
and the gate runs it over the workspace. Exceeding it fails the gate; the accepted fix is an
extraction, not a waiver, and there is no per-crate override. Three extractions are visible in the
tree as files that exist for no other reason:

\begin{itemize}
	\item \texttt{daemon/sys\_status.rs} and \texttt{store/sys\_status.rs} --- the same record from
	two sides: the daemon's view of system status and the store's persisted form, split when the
	health work of row 370 pushed a file over the line.
	\item \texttt{daemon/model\_roles.rs} --- the role-to-model mapping, extracted from the config
	parser.
	\item \texttt{nodes/coder/} --- the coder's runtime selection, output handling, skills and
	container path, each in its own file.
\end{itemize}

The second budget is conceptual and harder to enforce: one idea gets one implementation. A copy is
a defect even when both copies are correct today, because the only question that matters is what
happens the first time one of them is fixed. Two rows in this repository exist because the rule was
broken and then noticed:

\begin{itemize}
	\item Row 359, \enquote{delete the copies}: implementations that had been duplicated across
	crates were removed rather than kept in sync by hand.
	\item Row 371, \enquote{the idle daemon leaves no files}: the lint registry used to open scratch
	databases at literal paths, which failed outside one specific checkout and silently dropped node
	kinds from the registry --- so lint reported \enquote{unknown kind} for graphs that were fine.
	The fix moved the scratch stores into memory, which removed both the file and the failure mode.
\end{itemize}

That second one is a good example of the workspace's recurring shape: the bug was not in the rule,
it was in the copy of the environment the rule assumed. A lint that needs a writable directory is a
lint that behaves differently on a host where the directory is not writable, and the difference
shows up as a false refusal rather than as an error about the directory.

\section{The gate}

There is one step list, \texttt{gate-steps.txt}, and both the local gate and the CI-shaped gate
read it. Each line is a label and a shell command; a step passes only if the command exits zero,
and every step runs even after an earlier failure, so one run reports every problem rather than the
first. The list currently has twenty steps, among them:

\begin{center}
\footnotesize
\begin{tabular}{@{}p{0.22\textwidth}p{0.66\textwidth}@{}}
	\toprule
	\textbf{Step} & \textbf{What it refuses} \\
	\midrule
	\texttt{fmt} & any file not formatted by rustfmt \\
	\texttt{clippy} & any warning, with \texttt{-D warnings} \\
	\texttt{build} & a binary that does not link; later tests spawn the real \texttt{pii} \\
	\texttt{test (lib/bins/tests)}, \texttt{test (doc)} & a unit, integration or doctest failure \\
	\texttt{lint-files} & a file over 450 effective lines \\
	\texttt{lint-scratch} & scratch state written outside \texttt{\_tmp} \\
	\texttt{lint-observability} & an event name or metric outside the closed vocabulary \\
	\texttt{reachability-check} & a public item nothing reaches \\
	\texttt{deps-check} & a dependency outside the allowlist, or a forbidden edge \\
	\texttt{net-check} & a network call outside the declared surface \\
	\texttt{autonomy}, \texttt{plan-verify}, \texttt{item-ids} & prose drift in specs, plans and
	row ids \\
	\bottomrule
\end{tabular}
\end{center}

The gate has two modes that read the same list. The routine mode runs in the working tree, with a
warm cache, and is what a developer runs while working. The fresh mode exports \texttt{HEAD} with
\texttt{git archive} into a throwaway directory and runs the same steps there:

\begin{lstlisting}[language=bash]
cargo run -q -p xtask -- fresh-gate            # export HEAD, run gate-steps.txt in a clean tree
cargo run -q -p xtask -- fresh-gate --worktree # same, from the worktree instead of HEAD
bash quality-gate.sh                           # routine mode, warm cache, current tree
bash quality-gate.sh --summary                 # one line per step
\end{lstlisting}

The instruction that goes with it is blunt and was learned the hard way: when the two modes
disagree, believe the fresh one. A gate can be green in a tree that has an untracked file, a stale
target directory or a modified manifest, and red in a clean clone of the same commit. The fresh
mode is the release boundary; the routine mode is a convenience.

\begin{center}
\resizebox{\linewidth}{!}{%
\begin{tikzpicture}[node distance=8mm and 12mm]
	\node[b] (tree) {the working tree\\\scriptsize untracked files, warm target};
	\node[b, below=of tree] (head) {HEAD\\\scriptsize \texttt{git archive}};
	\node[b, right=of tree] (routine) {routine gate\\\scriptsize fast, same steps};
	\node[b, right=of head] (fresh) {fresh gate\\\scriptsize same steps, clean tree};
	\node[v, right=of routine] (agree) {agree?};
	\node[b, right=of fresh] (believe) {believe this one};

	\draw[f] (tree) -- (routine);
	\draw[f] (head) -- (fresh);
	\draw[f] (routine) -- (agree);
	\draw[sf] (fresh.north) to[bend left=10] node[p,above=2pt,midway]{compare} (agree.south);
	\draw[f] (agree) -- node[p,above=2pt,midway]{no} (believe);
\end{tikzpicture}}
\captionof{figure}{The two modes of one gate, which read the same step list on purpose, because a step added to one and
not the other would report a different truth depending on how it was invoked.}

\end{center}

\section{Four binaries}

The workspace produces four executables, and the split is by audience rather than by function.

\begin{itemize}
	\item \texttt{pii} (\texttt{pii-cli}) --- the operator's entry point, and the only one most
	users run. It embeds the daemon for commands that need one, and speaks to a running daemon over
	the unix socket for the rest.
	\item \texttt{pii-mcp} --- the MCP server, so an editor or another agent can call the same
	tools the turn loop calls.
	\item \texttt{pii-web} --- a small HTTP surface over the protocol, for reading runs and graphs
	from a browser.
	\item \texttt{pii-evolve} --- the driver for the evolution loop, which proposes changes from
	recorded trajectories.
\end{itemize}

Everything else is a library. The optional surfaces --- the TUI, the web viewer, the MCP server,
the WASM plugin host, the hardware abstraction --- sit behind their own features so that a build for
a host without GPIO does not compile GPIO code. That is also why the gate lints the real GPIO
backend without ever running it: code that is compiled but never exercised rots silently, and a
lint is the cheapest thing that can still be true about it.

\section{The coupling that is documented rather than removed}

One edge inside \texttt{pii-server} is worth naming because it explains a design decision that
looks unfinished. The JSON-RPC dispatch table in \texttt{engine/dispatch.rs} calls into
\texttt{engine/methods.rs} for roughly ten methods --- worker list and reload, approval list and
resolve, heal list and get, cron list, channel list, memory search/put/consolidate, system reload,
run resume. That is a real dependency from the engine-bridge concern into the daemon-application
concern, and it is the reason a full split into \texttt{pii-engine}, \texttt{pii-nodes} and
\texttt{pii-daemon} has not been done.

The alternative was a crate dependency edge that says the bridge depends on the application, which
is exactly the kind of edge the workspace spends its \texttt{deps-check} step preventing. The
decision recorded in the module header is to keep the coupling visible in one file, document it
there, and defer the split to a reviewed change rather than pretend the boundary is clean. That is
the general pattern in this workspace: where a rule cannot be satisfied yet, the violation is
written down where the next reader will trip over it, not smoothed over in a comment elsewhere.

\section{Finding a behaviour from a symptom}

The question a new contributor actually asks is not \enquote{which crate owns X} but \enquote{I
saw Y, where do I look}. The table is the answer, and every path in it is a real file.

\begin{longtable}{@{}p{0.36\textwidth}p{0.56\textwidth}@{}}
	\toprule
	\textbf{Symptom} & \textbf{Open first} \\
	\midrule
	\endhead
	A graph will not load, with a rule code & \texttt{crates/pii-graph/src/validate.rs} for
	\texttt{V*}; \texttt{fragment/call.rs} for \texttt{graph.call.*}; \texttt{fragment/splice.rs} for
	\texttt{fragment.*} \\
	A call refuses on capabilities & \texttt{fragment/call.rs}, the \texttt{widen} checks: one for
	the call's grant, one per callee node \\
	A queue row stays \texttt{Running} & \texttt{crates/pii-store/src/queue.rs} for the status
	transitions; \texttt{graphs/queue.toml} for the edges that decide whether a verdict ever reaches
	\texttt{record} \\
	A run never finishes and its journal stops mid-graph & \texttt{crates/pii-graph/src/run/engine.rs}
	(the tick loop) and \texttt{sched/state.rs} (\texttt{readiness}, \texttt{propagate\_dead}) \\
	A run resumed after a restart re-ran work & \texttt{run/resume.rs}: a node with no terminal
	event is re-run by design; look for the \texttt{NodeInterrupted} marker \\
	Health is slow & \texttt{engine/conn.rs} for the heavy-method set;
	\texttt{daemon/identity.rs} for the record it should be reading \\
	An agent left running after a crash & \texttt{engine/orphans.rs}, and the
	\texttt{coder.orphan\_reaped} event it journals \\
	A configuration key seems ignored & \texttt{daemon/config.rs}: every section is
	\texttt{deny\_unknown\_fields}, so an ignored key is a load error, not a default \\
	\bottomrule
\end{longtable}

Two habits make that table faster to use. The first is to grep for the rule code rather than the
prose: every refusal string is a literal that appears in exactly one place, so
\texttt{grep -rn "graph.call.widen" crates/} lands on the check. The second is to read the graph
file before the code. Most \enquote{engine} bugs in this repository have been edge bugs --- a port
with no edge, an edge to a node that can never be ready --- and the graph is a page long.

\begin{lstlisting}[language=bash]
grep -rn 'graph.call.widen' crates/            # which check raised it
grep -rn 'V1:' crates/pii-graph/src/validate.rs # the port rule and its message
pii graph lint graphs/queue.toml               # the same rules, over one file
pii graph show graphs/spec.toml                # nodes, ports and edges after expansion
pii run trace <run_id>                         # the journal, in sequence order
\end{lstlisting}

The third and fourth lines are the ones to reach for first when a graph is in question. \texttt{pii
graph lint} runs the loader and the static validator and prints every warning as well as every
error, and the warnings are where a missing \texttt{err} edge shows up: a declared port with no
outgoing edge is a warning, not an error, because a graph may legitimately stop on it. \texttt{pii
graph show} prints the graph after calls are expanded --- which is the only honest way to see what
will actually run for a graph that composes.
